% Generated by benchmarks/plot_sensitivity.py from validated summary.json.
\begin{table*}[t]
\centering
\caption{Sensitivity study: ranges across three interpreter sessions and the stated seeds. Warm speedups compare exact and bounded execution at equal worker counts. Lifetime charges eight collections, profiling, search and pool setup/shutdown. Payload is the largest individual complex128 state, not RSS. Budget: $0.01$ per output.}
\label{tab:sensitivity-results}
\footnotesize
\begin{tabular}{lrrrrrrr}
\toprule
Fixture & Seeds & Workers & Width & Payload factor & Warm speedup & Lifetime speedup & Profile (ms) \\
\midrule
8q / 2 blocks & 1 & 1 & $8\to 4$ & 16 & 2.23--2.40 & 0.66--0.76 & 7.50--8.07 \\
8q / 2 blocks & 1 & 2 & $8\to 4$ & 16 & 1.58--2.64 & 0.86--0.88 & 21.80--36.20 \\
8q / 2 blocks & 1 & 4 & $8\to 4$ & 16 & 1.16--1.79 & 0.82--0.85 & 25.95--28.15 \\
12q / 2 blocks & 1 & 1 & $12\to 6$ & 64 & 4.80--5.18 & 1.35--1.35 & 11.17--13.00 \\
12q / 2 blocks & 1 & 2 & $12\to 6$ & 64 & 2.48--3.66 & 0.93--0.97 & 21.71--24.60 \\
12q / 2 blocks & 1 & 4 & $12\to 6$ & 64 & 2.46--2.84 & 0.93--0.98 & 21.76--24.11 \\
16q / 2 blocks & 3 & 1 & $16\to 8$ & 256 & 40.56--53.88 & 2.89--3.57 & 133.05--165.46 \\
16q / 2 blocks & 3 & 2 & $16\to 8$ & 256 & 27.43--31.70 & 2.27--2.57 & 87.83--103.89 \\
16q / 2 blocks & 3 & 4 & $16\to 8$ & 256 & 22.70--30.28 & 2.04--2.49 & 89.91--109.28 \\
15q / 3 blocks & 3 & 1 & $15\to 5$ & 1024 & 30.95--35.46 & 1.31--1.55 & 138.32--158.77 \\
15q / 3 blocks & 3 & 2 & $15\to 5$ & 1024 & 9.31--12.51 & 0.88--1.02 & 138.80--181.29 \\
15q / 3 blocks & 3 & 4 & $15\to 5$ & 1024 & 11.86--17.78 & 0.99--1.10 & 106.85--137.25 \\
16q connected control & 1 & 1 & $16\to 16$ & 1 & 0.92--1.03 & 0.77--0.82 & 138.56--165.83 \\
16q connected control & 1 & 2 & $16\to 16$ & 1 & 0.97--1.00 & 0.83--0.85 & 106.49--130.78 \\
16q connected control & 1 & 4 & $16\to 16$ & 1 & 0.99--1.01 & 0.83--0.84 & 108.31--126.78 \\
\bottomrule
\end{tabular}
\end{table*}


\subsection{Proof of Theorem~\ref{thm:observable-error}}
\begin{IEEEproof}
For a Pauli rotation channel $\mathcal U_g$, define $\mathcal D_g(A)=-i[P_g/2,A]$. Then $\partial_g\mathcal U_g=\mathcal D_g\mathcal U_g$ and $\partial_g^2\mathcal U_g=\mathcal D_g^2\mathcal U_g$. Since $\|P_g\|_\infty=1$,
\begin{equation}
 \|\mathcal D_g(A)\|_1\leq\|P_gA\|_1/2+\|AP_g\|_1/2\leq\|A\|_1.
\end{equation}
Every unitary channel preserves trace norm. Differentiating a circuit once or twice inserts one or two such derivative maps. This also covers two derivatives of the same occurrence. Using $\|\rho\|_1=1$ and $|\operatorname{Tr}(O_k A)|\leq\|O_k\|_\infty\|A\|_1$ therefore gives
\begin{equation}
 |\partial_g f_k(\vartheta)|\leq1,\qquad
 |\partial_g\partial_h f_k(\vartheta)|\leq1
 \label{eq:uniform-derivative-bounds}
\end{equation}
at every angle vector $\vartheta$, not only the profiled point. A structural backward cone contains every gate that can affect the evolved observable. Changing angles leaves gate supports fixed, so derivatives involving an occurrence outside $C_k$ vanish identically.

Put $q(t)=f_k(\theta+t\delta)$. Equation~\eqref{eq:uniform-derivative-bounds} implies $|q'(t)|\leq L_k$ and $|q''(t)|\leq L_k^2$. Taylor's integral identity gives
\begin{equation}
 q(1)-q(0)=J_k\delta+\int_0^1(1-t)q''(t)\,dt.
\end{equation}
Taking absolute values bounds the remainder by $L_k^2/2$. Integrating the first-derivative bound yields $L_k$, while the range of a norm-one observable yields two. Taking their minimum proves Eq.~\eqref{eq:ideal-sensitivity-bound}.

For sharpness, compose $m$ single-qubit $R_Y$ occurrences on $|0\rangle$, all at zero, and request $Z$. Perturb every occurrence by $d$. The gradient at the starting point vanishes, but the output change is $1-\cos(md)$ with $L=m|d|$. Consequently,
\begin{equation}
 \lim_{d\to0}\frac{1-\cos(md)}{(m d)^2}=\frac12.
\end{equation}
No smaller uniform coefficient can replace $1/2$ in the quadratic remainder. The same example shows why counting shared occurrences is necessary. Expanding around a total rotation of $\pi/2$ similarly saturates the linear bound to first order.

Finally, with the structural zeros enforced, $|(\widehat J_k-J_k)\delta|\leq\eta L_k$. Adding this to the Taylor term and allowing error $\eta$ at each independently evaluated endpoint proves Eq.~\eqref{eq:sensitivity-bound}. This last step is conditional on the stated numerical error bounds and exact evaluation of the bound expression itself.
\end{IEEEproof}

The theorem is an application of unitary derivative bounds and Taylor's integral remainder, not an inference from a finite sample. The two corollaries add sufficient precision and reduced-evaluation conditions for this planner. The experiments validate particular implementations and timing hypotheses; they do not establish a universal speedup or the novelty of the underlying mathematical inequalities.

\subsection{Sensitivity Reproducibility}

The implementation assumes $\eta=10^{-10}$ in the study. Independent double-precision agreement does not prove this allowance or establish outward rounding of the computed bound. The bound is consequently an analytic error-control rule with an explicit numerical assumption. The evaluated program must already be feasible; neither small gradients nor approximate candidates supply otherwise unavailable exact profiles.

\path{src/lazyqc/sensitivity.py} implements profiling and bounded candidate search. \path{benchmarks/run_sensitivity.py} defines the cases; \path{tests/test_sensitivity.py} includes independent dense-matrix and analytic regression tests. Raw final observations reside in \path{results/sensitivity/session0/} through \path{results/sensitivity/session2/}; \path{results/sensitivity/frozen_source/} contains their measured source snapshot. Pilot runs are excluded from aggregation. \path{benchmarks/summarize_sensitivity.py} validates source hashes and produces the summary; \path{docs/SENSITIVITY.md} specifies cost scopes and reproduction. Native thread counts are fixed at one on the same unpinned, heterogeneous i7-13700H host used for the local-worker study. Summed RSS includes shared pages more than once and may retain allocations from earlier exact collections.
